\section{Summation of Primes --- Problem 10}
\label{prob:10}
\noindent\textbf{Source:} \url{https://github.com/KO0I/projectEuler/blob/reorganize/solutions-by-language/solutions/c/summation_of_primes_10.c}

\noindent\textbf{Assessment:} Computes a different sum.

\subsection*{Mathematical goal}
The problem asks for the sum of primes strictly below two million:
\[
S=\sum_{\substack{p<2\,000\,000\\p\text{ prime}}}p.
\]

\subsection*{How the C code works}
The program reads an integer \verb|n| from its command-line argument and loops from \verb|i=1| through \verb|n|, adding each integer to \verb|sum|. It thus computes the arithmetic series \(n(n+1)/2\). There is no primality test.

\subsection*{Correctness and limitations}
This C program does not solve the prime-summation problem. Using \verb|int| also overflows for the intended bound. A sieve of Eratosthenes combined with a 64-bit accumulator would solve the actual problem. The mathematical answer is 142913828922, not the sum produced by the present algorithm.

